\chapter{Section graded ring infrastructure: tensor powers and graded sections}
\label{chap:Picard_SectionGradedRing}

\providecommand{\PshMod}{\mathrm{PresheafOfModules}}
\providecommand{\ShMod}{\mathrm{SheafOfModules}}
\providecommand{\Modules}{\mathrm{Modules}}

The Hilbert-polynomial route of \cref{chap:Picard_QuotScheme} represents the
degree-\(m\) graded component of a fibre by the global sections
\(\Gamma(X_s, \mathcal{F}_s \otimes L_s^{\otimes m})\) of a twist of a coherent
sheaf by a tensor power of a line bundle, and it organises these components into
a graded ring \(R(X_s, L_s) = \bigoplus_{m \ge 0} \Gamma(X_s, L_s^{\otimes m})\)
(\cref{def:sectionGradedRing}) and a graded module
\(M(X_s, \mathcal{F}_s, L_s) = \bigoplus_{m \ge 0}
\Gamma(X_s, \mathcal{F}_s \otimes L_s^{\otimes m})\)
(\cref{def:sectionGradedModule}). Three pieces of that picture have no direct
Mathlib counterpart and must be built here, each in its own section:

\begin{enumerate}
  \item \textbf{Tensor powers of a sheaf of modules}
    (\cref{sec:sgr_tensor_powers}). Mathlib equips the category
    \(\PshMod\) of \emph{presheaves} of modules over a presheaf of commutative
    rings with a symmetric monoidal structure
    (\cref{lem:presheafModule_monoidal_mathlib}) and provides the
    sheafification functor \(\PshMod \to \ShMod\)
    (\cref{lem:presheafModule_sheafification_mathlib}); but the category
    \(\ShMod\) of \emph{sheaves} of modules carries \emph{no} monoidal
    structure. The tensor product of sheaves of modules
    \(\mathcal{F} \otimes_{\mathcal{O}_X} \mathcal{G}\) must therefore be built
    as the sheafification of the pointwise tensor presheaf, and the iterated
    tensor power \(\mathcal{L}^{\otimes m}\) recursively from it.
  \item \textbf{Lax-monoidal global sections} (\cref{sec:sgr_lax_sections}).
    The global-sections functor \(\Gamma\) is only \emph{lax} monoidal: from the
    pointwise-strict monoidality of the presheaf tensor product and the
    sheafification unit one obtains a natural multiplication
    \(\Gamma(\mathcal{F}) \otimes_{\Gamma(\mathcal{O}_X)} \Gamma(\mathcal{G})
    \to \Gamma(\mathcal{F} \otimes_{\mathcal{O}_X} \mathcal{G})\), together with
    the associativity and unit coherence those products must satisfy.
  \item \textbf{Graded-ring / graded-module assembly}
    (\cref{sec:sgr_graded_assembly}). The section multiplications and the
    tensor-power comparison isomorphisms assemble, through Mathlib's
    external-direct-sum graded API (\cref{lem:directSum_gcommSemiring_mathlib},
    \cref{lem:directSum_gmodule_mathlib}), into the graded \emph{ring} structure
    on \(\bigoplus_{m \ge 0} \Gamma(X_s, L_s^{\otimes m})\) and the graded
    \emph{module} structure on
    \(\bigoplus_{m \ge 0} \Gamma(X_s, \mathcal{F}_s \otimes L_s^{\otimes m})\).
\end{enumerate}

\section{Tensor powers of a sheaf of modules}
\label{sec:sgr_tensor_powers}

We work over a fixed scheme \(X\) (in the application \(X = X_s\), a fibre over a
field \(\kappa(s)\)). Its category of sheaves of modules is
\(X.\Modules = \ShMod(\mathcal{O}_X)\), with global-sections notation
\(\Gamma(X, \mathcal{M}) = \mathcal{M}(X)\) for the value of the underlying
presheaf on the top open. The structure sheaf \(\mathcal{O}_X\) is a sheaf of
commutative rings, so its underlying presheaf factors through the category of
commutative rings; this is the hypothesis under which Mathlib's presheaf-level
monoidal structure applies.

\begin{lemma}[Monoidal structure on presheaves of modules]
  \label{lem:presheafModule_monoidal_mathlib}
  \lean{PresheafOfModules.monoidalCategory}
  \mathlibok
  \textit{Provided by Mathlib
  (\texttt{Mathlib.Algebra.Category.ModuleCat.Presheaf.Monoidal}).}
  Let \(R : \mathcal{C}^{\mathrm{op}} \to \mathrm{CommRingCat}\) be a presheaf of
  commutative rings. Then the category \(\PshMod(R)\) of presheaves of modules
  over \(R\) carries a symmetric monoidal structure whose tensor product is
  computed objectwise: for presheaves of modules \(\mathcal{M}_1, \mathcal{M}_2\)
  and an object \(U\),
  \[
    (\mathcal{M}_1 \otimes \mathcal{M}_2)(U)
      \;=\; \mathcal{M}_1(U) \otimes_{R(U)} \mathcal{M}_2(U),
  \]
  with restriction maps the tensor products of the restriction maps. The unit is
  the presheaf \(R\) viewed as a module over itself, and the associator,
  unitors, and braiding are induced objectwise from those of
  \(\mathrm{ModuleCat}\). In particular evaluation at a fixed object \(U\) is a
  \emph{strict} monoidal functor
  \(\PshMod(R) \to \mathrm{ModuleCat}(R(U))\).
\end{lemma}

\begin{lemma}[Sheafification of presheaves of modules]
  \label{lem:presheafModule_sheafification_mathlib}
  \lean{PresheafOfModules.sheafification}
  \mathlibok
  \textit{Provided by Mathlib
  (\texttt{Mathlib.Algebra.Category.ModuleCat.Presheaf.Sheafification}).}
  Let \(R\) be a sheaf of (commutative) rings on a site, with underlying
  presheaf \(R_0\). There is a sheafification functor
  \[
    (-)^{\#} \;:\; \PshMod(R_0) \longrightarrow \ShMod(R),
  \]
  left adjoint to the forgetful functor, together with the unit
  \(\eta_{\mathcal{M}_0} : \mathcal{M}_0 \to (\mathcal{M}_0)^{\#}\) of the
  adjunction. On underlying sheaves of abelian groups it agrees with the
  ordinary sheafification, and it sends a presheaf of \(R_0\)-modules to the
  associated sheaf of \(R\)-modules.
\end{lemma}

\begin{lemma}[The structure sheaf as the unit module]
  \label{lem:moduleUnit_mathlib}
  \lean{SheafOfModules.unit}
  \mathlibok
  \textit{Provided by Mathlib
  (\texttt{Mathlib.Algebra.Category.ModuleCat.Sheaf}).}
  For a sheaf of rings \(R\) (in the application \(R = \mathcal{O}_X\)) the
  structure sheaf, viewed as a module over itself, is a distinguished object
  \(\mathbf{1} = \mathcal{O}_X \in \ShMod(\mathcal{O}_X)\). It is the value at
  the zeroth tensor power, \(\mathcal{L}^{\otimes 0} = \mathcal{O}_X\), and its
  global sections \(\Gamma(X, \mathcal{O}_X)\) form the degree-zero component of
  the section graded ring.
\end{lemma}

\begin{definition}\leanok
[Presheaf-of-modules category in monoidal form]
  \label{def:monoidalPresheaf}
  \lean{AlgebraicGeometry.Scheme.Modules.MonoidalPresheaf}
  \uses{lem:presheafModule_monoidal_mathlib}
  For a scheme \(X\), write \(\PshMod_{\otimes}(X)\) for the category of presheaves
  of modules over the underlying presheaf of (plain) rings of \(\mathcal{O}_X\),
  obtained by composing the structure sheaf with the forgetful functor
  \(\mathrm{CommRingCat} \to \mathrm{RingCat}\). This is definitionally
  \(X.\PshMod\), but presented in the exact form for which Mathlib equips it with
  the symmetric monoidal structure of \cref{lem:presheafModule_monoidal_mathlib};
  it is the carrier on which the objectwise tensor product of
  \cref{def:sheafTensorObj} is taken.
\end{definition}

\begin{definition}\leanok
[Scheme-level module sheafification functor]
  \label{def:schemeModuleSheafification}
  \lean{AlgebraicGeometry.Scheme.Modules.sheafification}
  \uses{lem:presheafModule_sheafification_mathlib}
  For a scheme \(X\), the \emph{module sheafification functor}
  \((-)^{\#} : X.\PshMod \to X.\Modules\) is Mathlib's sheafification functor
  (\cref{lem:presheafModule_sheafification_mathlib}) specialised to the identity
  morphism of the underlying presheaf of rings of \(\mathcal{O}_X\) (a locally
  bijective map). Applying it to the objectwise tensor presheaf is what defines the
  sheaf tensor product \cref{def:sheafTensorObj}.
\end{definition}

\begin{definition}\leanok
[Unit module of a scheme]
  \label{def:unitModule}
  \lean{AlgebraicGeometry.Scheme.Modules.unitModule}
  \uses{lem:moduleUnit_mathlib}
  For a scheme \(X\), the \emph{unit module} \(\mathbf{1}_X \in X.\Modules\) is the
  structure sheaf viewed as a module over itself, i.e.\ Mathlib's
  \(\mathrm{SheafOfModules.unit}\) (\cref{lem:moduleUnit_mathlib}) for
  \(\mathcal{O}_X\). It is the value of the zeroth tensor power,
  \(\mathcal{L}^{\otimes 0} = \mathbf{1}_X\) (\cref{def:sheafTensorPow}).
\end{definition}

\begin{definition}\leanok
[Tensor product of sheaves of modules]
  \label{def:sheafTensorObj}
  \lean{AlgebraicGeometry.Scheme.Modules.tensorObj}
  \uses{lem:presheafModule_monoidal_mathlib,
    lem:presheafModule_sheafification_mathlib}
  % SOURCE: [Stacks Project], "Sheaves of Modules", \S "Tensor product" (Tag
  % 01CA), prose definition (read from references/stacks-modules.tex,
  % L2282-L2297).
  % SOURCE QUOTE:
  % "Let $(X, \mathcal{O}_X)$ be a ringed space and let
  %  $\mathcal{F}$ and $\mathcal{G}$ be $\mathcal{O}_X$-modules.
  %  We define first the {\it tensor product presheaf}
  %  $$
  %  \mathcal{F} \otimes_{p, \mathcal{O}_X} \mathcal{G}
  %  $$
  %  as the rule which assigns to $U \subset X$ open the
  %  $\mathcal{O}_X(U)$-module
  %  $\mathcal{F}(U) \otimes_{\mathcal{O}_X(U)} \mathcal{G}(U)$.
  %  Having defined this we define the {\it tensor product sheaf}
  %  as the sheafification of the above:
  %  $$
  %  \mathcal{F} \otimes_{\mathcal{O}_X} \mathcal{G}
  %  =
  %  (\mathcal{F} \otimes_{p, \mathcal{O}_X} \mathcal{G})^\#
  %  $$"
  \textit{Source: [Stacks Project], "Sheaves of Modules", Tag 01CA (tensor
  product).}
  Let \(\mathcal{F}, \mathcal{G}\) be sheaves of \(\mathcal{O}_X\)-modules. Their
  \emph{tensor product} is the sheaf of \(\mathcal{O}_X\)-modules
  \[
    \mathcal{F} \otimes_{\mathcal{O}_X} \mathcal{G}
      \;=\; \bigl(\mathcal{F} \otimes_{p,\mathcal{O}_X} \mathcal{G}\bigr)^{\#},
  \]
  the sheafification (\cref{lem:presheafModule_sheafification_mathlib}) of the
  tensor product presheaf \(\mathcal{F} \otimes_{p,\mathcal{O}_X} \mathcal{G}\),
  whose value on an open \(U\) is the \(\mathcal{O}_X(U)\)-module
  \(\mathcal{F}(U) \otimes_{\mathcal{O}_X(U)} \mathcal{G}(U)\). Concretely, the
  tensor product presheaf is the Mathlib presheaf-level monoidal product
  (\cref{lem:presheafModule_monoidal_mathlib}) of the underlying presheaves of
  modules of \(\mathcal{F}\) and \(\mathcal{G}\), and
  \(\mathcal{F} \otimes_{\mathcal{O}_X} \mathcal{G}\) is its sheafification. This
  construction is symmetric and associative because the presheaf tensor product
  is, and sheafification preserves the coherence isomorphisms.
\end{definition}

\begin{definition}\leanok
[Tensor power of a line bundle]
  \label{def:sheafTensorPow}
  \lean{AlgebraicGeometry.Scheme.Modules.tensorPow}
  \uses{def:sheafTensorObj, def:unitModule, lem:moduleUnit_mathlib}
  % SOURCE: [Stacks Project], "Sheaves of Modules", \S "Invertible modules",
  % `definition-powers' (Tag 01CU) (read from references/stacks-modules.tex,
  % L4220-L4228).
  % SOURCE QUOTE:
  % "Let $(X, \mathcal{O}_X)$ be a ringed space. Given an invertible sheaf
  %  $\mathcal{L}$ on $X$ and $n \in \mathbf{Z}$ we define the
  %  $n$th {\it tensor power} $\mathcal{L}^{\otimes n}$ of $\mathcal{L}$
  %  as the image of $\mathcal{O}_X$ under applying the equivalence
  %  $\mathcal{F} \mapsto\mathcal{F} \otimes_{\mathcal{O}_X} \mathcal{L}$
  %  exactly $n$ times."
  \textit{Source: [Stacks Project], "Sheaves of Modules", Tag 01CU (tensor
  powers).}
  Let \(\mathcal{L}\) be a sheaf of \(\mathcal{O}_X\)-modules (in the application
  an invertible sheaf / line bundle \(L_s\)). For \(m \in \mathbb{N}\) the
  \emph{\(m\)th tensor power} \(\mathcal{L}^{\otimes m}\) is defined by recursion
  on \(m\):
  \[
    \mathcal{L}^{\otimes 0} = \mathcal{O}_X
      \quad(\text{the unit module, \cref{lem:moduleUnit_mathlib}}),
    \qquad
    \mathcal{L}^{\otimes (m+1)}
      = \mathcal{L}^{\otimes m} \otimes_{\mathcal{O}_X} \mathcal{L}
        \quad(\cref{def:sheafTensorObj}).
  \]
  We restrict to non-negative powers, which is all the section graded ring uses;
  the Stacks definition extends \(\mathcal{L}^{\otimes n}\) to all
  \(n \in \mathbf{Z}\) using invertibility, but negative powers are not needed
  here. For \(\mathcal{L}\) invertible each \(\mathcal{L}^{\otimes m}\) is again
  invertible.
\end{definition}

\begin{lemma}\leanok
[Zeroth tensor power is the unit]
  \label{lem:tensorPow_zero}
  \lean{AlgebraicGeometry.Scheme.Modules.tensorPow_zero}
  \uses{def:sheafTensorPow, def:unitModule}
  For a sheaf of modules \(\mathcal{L}\) on \(X\) one has
  \(\mathcal{L}^{\otimes 0} = \mathbf{1}_X\), the unit module
  (\cref{def:unitModule}); this is the base clause of the recursive definition
  \cref{def:sheafTensorPow}, holding by definition.
\end{lemma}

\begin{lemma}\leanok
[Tensor-power recursion step]
  \label{lem:tensorPow_succ}
  \lean{AlgebraicGeometry.Scheme.Modules.tensorPow_succ}
  \uses{def:sheafTensorPow, def:sheafTensorObj}
  For a sheaf of modules \(\mathcal{L}\) on \(X\) and \(m \in \mathbb{N}\) one has
  \(\mathcal{L}^{\otimes (m+1)} = \mathcal{L}^{\otimes m} \otimes_{\mathcal{O}_X}
  \mathcal{L}\) (\cref{def:sheafTensorObj}); this is the successor clause of
  \cref{def:sheafTensorPow}, holding by definition.
\end{lemma}

\begin{definition}\leanok
[Twist of a coherent sheaf by a tensor power]
  \label{def:sheafModuleTwist}
  \lean{AlgebraicGeometry.Scheme.Modules.moduleTensorPow}
  \uses{def:sheafTensorObj, def:sheafTensorPow}
  Let \(\mathcal{F}\) be a sheaf of \(\mathcal{O}_X\)-modules (in the application
  a coherent sheaf \(\mathcal{F}_s\)) and \(\mathcal{L}\) a line bundle. For
  \(m \in \mathbb{N}\) the \emph{\(m\)-twist} of \(\mathcal{F}\) by
  \(\mathcal{L}\) is
  \[
    \mathcal{F}(m) \;:=\;
      \mathcal{F} \otimes_{\mathcal{O}_X} \mathcal{L}^{\otimes m}
      \qquad(\cref{def:sheafTensorObj}, \cref{def:sheafTensorPow}),
  \]
  the degree-\(m\) carrier of the section graded module. By construction
  \(\mathcal{F}(0) = \mathcal{F} \otimes_{\mathcal{O}_X} \mathcal{O}_X \cong
  \mathcal{F}\).
\end{definition}

\begin{lemma}\leanok
[Zeroth twist]
  \label{lem:moduleTensorPow_zero}
  \lean{AlgebraicGeometry.Scheme.Modules.moduleTensorPow_zero}
  \uses{def:sheafModuleTwist, def:sheafTensorObj, def:unitModule,
    def:tensorObjRightUnitor}
  For sheaves of modules \(\mathcal{F}, \mathcal{L}\) on \(X\) one has
  \(\mathcal{F}(0) = \mathcal{F} \otimes_{\mathcal{O}_X} \mathbf{1}_X\)
  (\cref{def:sheafTensorObj}, \cref{def:unitModule}); this is the base clause of
  \cref{def:sheafModuleTwist}, holding by definition, and is canonically
  isomorphic to \(\mathcal{F}\) via the right unitor
  (\cref{def:tensorObjRightUnitor}).
\end{lemma}

The next four isomorphisms are the parts of the (would-be) symmetric monoidal
structure on \(X.\Modules\) that descend through sheafification from
\cref{lem:presheafModule_monoidal_mathlib} using only \emph{functoriality} of the
sheafification functor (\cref{def:schemeModuleSheafification}) and, for the
unitors, the reflective counit isomorphism --- no strong monoidality of
sheafification is required. They are the structural inputs of the tensor-power
comparison \cref{lem:sheafTensorPow_add}.

\begin{definition}\leanok
[Counit isomorphism of module sheafification]
  \label{def:sheafificationCounitIso}
  \lean{AlgebraicGeometry.Scheme.Modules.sheafificationCounitIso}
  \uses{def:schemeModuleSheafification, lem:presheafModule_sheafification_mathlib}
  For a sheaf of modules \(\mathcal{G}\) on \(X\) the counit of the sheafification
  adjunction (\cref{lem:presheafModule_sheafification_mathlib},
  \cref{def:schemeModuleSheafification}) gives a canonical isomorphism
  \((\mathcal{G})^{\#} \xrightarrow{\ \sim\ } \mathcal{G}\) of sheaves of modules,
  invertible because the right adjoint forgetful functor to presheaves of modules
  is fully faithful.
\end{definition}

\begin{definition}\leanok
[Left unitor of the sheaf tensor product]
  \label{def:tensorObjUnitIso}
  \lean{AlgebraicGeometry.Scheme.Modules.tensorObjUnitIso}
  \uses{def:sheafTensorObj, def:unitModule, def:sheafificationCounitIso,
    lem:presheafModule_monoidal_mathlib}
  For a sheaf of modules \(\mathcal{G}\) on \(X\) there is a canonical left-unitor
  isomorphism \(\mathbf{1}_X \otimes_{\mathcal{O}_X} \mathcal{G}
  \xrightarrow{\ \sim\ } \mathcal{G}\) (\cref{def:sheafTensorObj},
  \cref{def:unitModule}), obtained by sheafifying the presheaf left unitor of
  \cref{lem:presheafModule_monoidal_mathlib} and composing with the counit
  isomorphism \cref{def:sheafificationCounitIso}. It is the base case
  (\(m = 0\)) of \cref{lem:sheafTensorPow_add}.
\end{definition}

\begin{definition}\leanok
[Right unitor of the sheaf tensor product]
  \label{def:tensorObjRightUnitor}
  \lean{AlgebraicGeometry.Scheme.Modules.tensorObjRightUnitor}
  \uses{def:sheafTensorObj, def:unitModule, def:sheafificationCounitIso,
    lem:presheafModule_monoidal_mathlib}
  For a sheaf of modules \(\mathcal{G}\) on \(X\) there is a canonical right-unitor
  isomorphism \(\mathcal{G} \otimes_{\mathcal{O}_X} \mathbf{1}_X
  \xrightarrow{\ \sim\ } \mathcal{G}\) (\cref{def:sheafTensorObj},
  \cref{def:unitModule}), the sheafification of the presheaf right unitor
  (\cref{lem:presheafModule_monoidal_mathlib}) composed with the counit
  isomorphism \cref{def:sheafificationCounitIso}.
\end{definition}

\begin{definition}\leanok
[Braiding of the sheaf tensor product]
  \label{def:tensorBraiding}
  \lean{AlgebraicGeometry.Scheme.Modules.tensorBraiding}
  \uses{def:sheafTensorObj, def:schemeModuleSheafification,
    lem:presheafModule_monoidal_mathlib}
  For sheaves of modules \(\mathcal{F}, \mathcal{G}\) on \(X\) there is a canonical
  braiding isomorphism \(\mathcal{F} \otimes_{\mathcal{O}_X} \mathcal{G}
  \xrightarrow{\ \sim\ } \mathcal{G} \otimes_{\mathcal{O}_X} \mathcal{F}\)
  (\cref{def:sheafTensorObj}), obtained by sheafifying
  (\cref{def:schemeModuleSheafification}) the symmetric braiding of the presheaf
  monoidal structure (\cref{lem:presheafModule_monoidal_mathlib}). It is the
  symmetry used in the inductive step of \cref{lem:sheafTensorPow_add}.
\end{definition}

The strong-monoidality comparison \cref{lem:isIso_sheafification_whiskerRight_unit}
below is assembled from three reduction lemmas --- a localization criterion, the
local-isomorphism property of the unit, and their corollary --- together with the
coequalizer presentation of the relative module-presheaf tensor product, the single
Mathlib-absent brick on which the comparison rests.

\begin{definition}

\leanok
[Grothendieck topology on the opens of a scheme]
  \label{def:snap_opensTopology}
  \lean{AlgebraicGeometry.Scheme.Modules.opensTopology}
  For a scheme \(X\) the \emph{opens topology} \(J\) is the canonical Grothendieck
  topology on the poset \((\mathrm{Opens}\,X)\) of opens of the underlying
  topological space, in which a sieve covers an open \(U\) precisely when the opens
  it contains cover \(U\). It is the site on which the abelian sheafification, and
  the weak-equivalence class \(J.W\) of local isomorphisms of abelian-group
  presheaves used in \cref{lem:isIso_sheafification_map_iff} and
  \cref{lem:localIso_toPresheaf_map_unit}, are taken.
\end{definition}

\begin{lemma}\leanok
[Localization criterion for module sheafification]
  \label{lem:isIso_sheafification_map_iff}
  \lean{AlgebraicGeometry.Scheme.Modules.isIso_sheafification_map_iff}
  \uses{def:schemeModuleSheafification, lem:presheafModule_sheafification_mathlib,
    def:snap_opensTopology}
  For a morphism \(f : \mathcal{P} \to \mathcal{Q}\) of presheaves of
  \(\mathcal{O}_X\)-modules, the sheafification \((f)^{\#}\)
  (\cref{def:schemeModuleSheafification}) is an isomorphism of sheaves of modules if
  and only if the underlying morphism of abelian-group presheaves lies in the
  weak-equivalence class \(J.W\) of the opens topology \(J\) on \(X\) (i.e.\ is a
  local isomorphism of abelian-group presheaves).
\end{lemma}

\begin{proof}
  \uses{lem:presheafModule_sheafification_mathlib}
  \leanok
  Module sheafification (\cref{def:schemeModuleSheafification}) is the localization
  functor at the class \(W\) of morphisms of presheaves of modules whose underlying
  abelian-presheaf morphism lies in \(J.W\), i.e.\ the inverse image of \(J.W\) along
  the forgetful functor. The localization criterion for a reflective adjunction
  (\cref{lem:presheafModule_sheafification_mathlib}) then gives the stated
  equivalence.
\end{proof}

\begin{lemma}\leanok
[The sheafification unit is an abelian local isomorphism]
  \label{lem:localIso_toPresheaf_map_unit}
  \lean{AlgebraicGeometry.Scheme.Modules.localIso_toPresheaf_map_unit}
  \uses{def:schemeModuleSheafification, lem:presheafModule_sheafification_mathlib,
    def:snap_opensTopology}
  For a presheaf of \(\mathcal{O}_X\)-modules \(\mathcal{P}\) the underlying
  abelian-group presheaf morphism of the sheafification unit
  \(\eta_{\mathcal{P}} : \mathcal{P} \to (\mathcal{P})^{\#}\)
  (\cref{lem:presheafModule_sheafification_mathlib}) is the abelian sheafification
  unit \(\mathcal{P} \to a(\mathcal{P})\); it therefore lies in the weak-equivalence
  class \(J.W\) of the opens topology.
\end{lemma}

\begin{proof}
  \uses{lem:presheafModule_sheafification_mathlib}
  \leanok
  The forgetful functor to abelian-group presheaves carries the module
  sheafification unit to the unit of the abelian associated-sheaf adjunction, which
  is a local isomorphism, hence a member of \(J.W\).
\end{proof}

\begin{lemma}\leanok
[Sheafification inverts the localization unit]
  \label{lem:isIso_sheafification_map_unit}
  \lean{AlgebraicGeometry.Scheme.Modules.isIso_sheafification_map_unit}
  \uses{lem:isIso_sheafification_map_iff, lem:localIso_toPresheaf_map_unit}
  For a presheaf of \(\mathcal{O}_X\)-modules \(\mathcal{P}\) the sheafification
  \((\eta_{\mathcal{P}})^{\#}\) of the localization unit is an isomorphism of sheaves
  of modules. This is the un-whiskered (\(m = 0\)) special case of
  \cref{lem:isIso_sheafification_whiskerRight_unit}.
\end{lemma}

\begin{proof}
  \uses{lem:isIso_sheafification_map_iff, lem:localIso_toPresheaf_map_unit}
  \leanok
  Feed the abelian local-isomorphism statement
  \cref{lem:localIso_toPresheaf_map_unit} through the localization criterion
  \cref{lem:isIso_sheafification_map_iff}.
\end{proof}

\begin{definition}

\leanok
[Abelian-carrier restriction map of a presheaf of modules]
  \label{def:snap_objRestrict}
  \lean{AlgebraicGeometry.Scheme.Modules.objRestrict}
  For a presheaf of \(\mathcal{O}_X\)-modules \(\mathcal{P}\) and a morphism
  \(f : U \to V\) in \((\mathrm{Opens}\,X)^{\mathrm{op}}\) (an inclusion of opens
  \(V \subseteq U\)), the \emph{abelian-carrier restriction map}
  \[
    \mathrm{objRestrict}_{\mathcal{P}}(f) :
      \mathcal{P}(U) \longrightarrow \mathcal{P}(V)
  \]
  is the underlying additive map of the presheaf restriction \(\mathcal{P}(f)\),
  viewed as a \(\mathbb{Z}\)-linear map between the carrier abelian groups. It is
  the form in which the presheaf restriction enters the \(\mathbb{Z}\)-tensor
  presheaves \cref{def:relTensorDomainPresheaf} and
  \cref{def:relTensorTriplePresheaf}; its carriers are pinned to print as
  \(\mathcal{P}(U), \mathcal{P}(V)\) so that the tensor-of-restrictions maps unify
  on elementary tensors.
\end{definition}

\begin{lemma}

\leanok
[Abelian-carrier restriction on elements]
  \label{lem:snap_objRestrict_apply}
  \lean{AlgebraicGeometry.Scheme.Modules.objRestrict_apply}
  \uses{def:snap_objRestrict}
  For a presheaf of modules \(\mathcal{P}\), a morphism \(f\) of opens, and a
  section \(x \in \mathcal{P}(U)\), one has
  \(\mathrm{objRestrict}_{\mathcal{P}}(f)\,(x) = \mathcal{P}(f)\,(x)\): the
  abelian-carrier restriction agrees with the presheaf restriction on elements, by
  definition.
\end{lemma}

\begin{lemma}

\leanok
[Identity law for the abelian-carrier restriction]
  \label{lem:snap_objRestrict_id}
  \lean{AlgebraicGeometry.Scheme.Modules.objRestrict_id}
  \uses{def:snap_objRestrict}
  For a presheaf of modules \(\mathcal{P}\) and an open \(U\),
  \(\mathrm{objRestrict}_{\mathcal{P}}(\mathrm{id}_U) = \mathrm{id}\): restriction
  along the identity is the identity, since \(\mathcal{P}\) preserves identities.
\end{lemma}

\begin{lemma}

\leanok
[Composition law for the abelian-carrier restriction]
  \label{lem:snap_objRestrict_comp}
  \lean{AlgebraicGeometry.Scheme.Modules.objRestrict_comp}
  \uses{def:snap_objRestrict}
  For a presheaf of modules \(\mathcal{P}\) and composable morphisms \(f : U \to V\),
  \(g : V \to W\) of opens, restriction along the composite is the composite of the
  restrictions,
  \[
    \mathrm{objRestrict}_{\mathcal{P}}(f \,;\, g)
      = \mathrm{objRestrict}_{\mathcal{P}}(g) \circ
        \mathrm{objRestrict}_{\mathcal{P}}(f),
  \]
  since \(\mathcal{P}\) preserves composition.
\end{lemma}

\begin{definition}\leanok
[Step-1 objectwise \(\mathbb{Z}\)-tensor presheaf]
  \label{def:relTensorDomainPresheaf}
  \lean{AlgebraicGeometry.Scheme.Modules.relTensorDomainPresheaf}
  \uses{lem:presheafModule_monoidal_mathlib, def:snap_objRestrict,
    lem:snap_objRestrict_id, lem:snap_objRestrict_comp}
  For presheaves of \(\mathcal{O}_X\)-modules \(\mathcal{P}, \mathcal{Q}\) the
  assignment \(U \mapsto \mathcal{P}(U) \otimes_{\mathbb{Z}} \mathcal{Q}(U)\),
  with restriction maps the \(\mathbb{Z}\)-tensor of the underlying restriction
  maps, is a functor \((\mathrm{Opens}\,X)^{\mathrm{op}} \to \mathrm{Ab}\). It is
  the codomain row of the relative-tensor coequalizer presentation
  (\cref{lem:relativeTensor_as_coequalizer}); functoriality follows from the
  functoriality of \(\mathrm{TensorProduct.map}\) (identity and composition laws
  by \(\otimes\)-induction).
\end{definition}

\begin{definition}\leanok
[Step-1 objectwise \(\mathbb{Z}\)-tensor triple presheaf]
  \label{def:relTensorTriplePresheaf}
  \lean{AlgebraicGeometry.Scheme.Modules.relTensorTriplePresheaf}
  \uses{lem:presheafModule_monoidal_mathlib, def:snap_objRestrict,
    lem:snap_objRestrict_id, lem:snap_objRestrict_comp}
  For presheaves of \(\mathcal{O}_X\)-modules \(\mathcal{P}, \mathcal{Q}\) the
  assignment
  \[
    U \;\longmapsto\;
      \mathcal{P}(U) \otimes_{\mathbb{Z}}
        \bigl(\mathcal{O}_X(U) \otimes_{\mathbb{Z}} \mathcal{Q}(U)\bigr),
  \]
  with restriction maps the \(\mathbb{Z}\)-tensor of the three underlying
  restriction maps --- those of \(\mathcal{P}\), of the structure presheaf
  \(\mathcal{O}_X\) (the middle factor restricting along the ring restriction
  map), and of \(\mathcal{Q}\) --- is a functor
  \((\mathrm{Opens}\,X)^{\mathrm{op}} \to \mathrm{Ab}\). It is the \emph{domain}
  row of the relative-tensor coequalizer presentation
  (\cref{lem:relativeTensor_as_coequalizer}): objectwise it is exactly the
  triple tensor \(\mathcal{P}(U) \otimes_{\mathbb{Z}} \mathcal{O}_X(U)
  \otimes_{\mathbb{Z}} \mathcal{Q}(U)\) on which the two
  \(\mathcal{O}_X(U)\)-action maps of
  \cref{lem:relativeTensor_objectwise_coequalizer} act, its codomain row being
  the \(\mathbb{Z}\)-tensor presheaf \cref{def:relTensorDomainPresheaf}.
  Functoriality (the identity and composition laws of the restriction maps)
  follows from the functoriality of \(\mathrm{TensorProduct.map}\) --- the
  \(\mathrm{map\_id}\) and \(\mathrm{map\_comp}\) laws --- checked on elementary
  tensors by \(\otimes\)-induction.
\end{definition}

\begin{definition}\leanok
[Left-action natural transformation of the coequalizer]
  \label{def:relTensorActL}
  \lean{AlgebraicGeometry.Scheme.Modules.relTensorActL}
  \uses{def:relTensorTriplePresheaf, def:relTensorDomainPresheaf,
    lem:snap_objRestrict_apply}
  For presheaves of \(\mathcal{O}_X\)-modules \(\mathcal{P}, \mathcal{Q}\) the
  \emph{left-action} natural transformation
  \[
    a_L : \bigl(U \mapsto \mathcal{P}(U) \otimes_{\mathbb{Z}}
      (\mathcal{O}_X(U) \otimes_{\mathbb{Z}} \mathcal{Q}(U))\bigr)
    \;\longrightarrow\;
    \bigl(U \mapsto \mathcal{P}(U) \otimes_{\mathbb{Z}} \mathcal{Q}(U)\bigr)
  \]
  from the triple presheaf (\cref{def:relTensorTriplePresheaf}) to the
  \(\mathbb{Z}\)-tensor presheaf (\cref{def:relTensorDomainPresheaf}) is, on an
  open \(U\), the additive map
  \[
    \mathcal{P}(U) \otimes_{\mathbb{Z}}
      \bigl(\mathcal{O}_X(U) \otimes_{\mathbb{Z}} \mathcal{Q}(U)\bigr)
    \longrightarrow
    \mathcal{P}(U) \otimes_{\mathbb{Z}} \mathcal{Q}(U),
    \qquad
    m \otimes (s \otimes n) \longmapsto (s \cdot m) \otimes n,
  \]
  collapsing the middle ring factor through the scalar action of
  \(\mathcal{O}_X(U)\) on \(\mathcal{P}(U)\). This is exactly the objectwise
  action map \(\mathrm{actLmap}\) of
  \cref{lem:relativeTensor_objectwise_coequalizer}; together with the parallel
  right-action transformation it is one of the two rows whose coequalizer
  presents the relative tensor presheaf \(\mathcal{P} \otimes_p \mathcal{Q}\)
  (\cref{lem:relativeTensor_as_coequalizer}). Naturality in \(U\) is the
  compatibility of the module action with the restriction maps: the restriction
  of \(s \cdot m\) equals \((s|_V) \cdot (m|_V)\), so \(a_L\) commutes with the
  triple-tensor and \(\mathbb{Z}\)-tensor restriction maps; this is checked on
  elementary tensors by \(\otimes\)-induction.
\end{definition}

\begin{definition}\leanok
[Right-action natural transformation of the coequalizer]
  \label{def:relTensorActR}
  \lean{AlgebraicGeometry.Scheme.Modules.relTensorActR}
  \uses{def:relTensorTriplePresheaf, def:relTensorDomainPresheaf,
    lem:snap_objRestrict_apply}
  For presheaves of \(\mathcal{O}_X\)-modules \(\mathcal{P}, \mathcal{Q}\) the
  \emph{right-action} natural transformation
  \[
    a_R : \bigl(U \mapsto \mathcal{P}(U) \otimes_{\mathbb{Z}}
      (\mathcal{O}_X(U) \otimes_{\mathbb{Z}} \mathcal{Q}(U))\bigr)
    \;\longrightarrow\;
    \bigl(U \mapsto \mathcal{P}(U) \otimes_{\mathbb{Z}} \mathcal{Q}(U)\bigr)
  \]
  from the triple presheaf (\cref{def:relTensorTriplePresheaf}) to the
  \(\mathbb{Z}\)-tensor presheaf (\cref{def:relTensorDomainPresheaf}) is, on an
  open \(U\), the additive map
  \[
    \mathcal{P}(U) \otimes_{\mathbb{Z}}
      \bigl(\mathcal{O}_X(U) \otimes_{\mathbb{Z}} \mathcal{Q}(U)\bigr)
    \longrightarrow
    \mathcal{P}(U) \otimes_{\mathbb{Z}} \mathcal{Q}(U),
    \qquad
    m \otimes (s \otimes n) \longmapsto m \otimes (s \cdot n),
  \]
  collapsing the middle ring factor through the scalar action of
  \(\mathcal{O}_X(U)\) on \(\mathcal{Q}(U)\). This is exactly the objectwise
  action map \(\mathrm{actRmap}\) of
  \cref{lem:relativeTensor_objectwise_coequalizer}, the mirror of the left-action
  transformation \(a_L\) (\cref{def:relTensorActL}); together \(a_L\) and \(a_R\)
  form the parallel pair whose coequalizer presents the relative tensor presheaf
  \(\mathcal{P} \otimes_p \mathcal{Q}\). Naturality in \(U\) is the compatibility
  of the module action with the restriction maps: the restriction of
  \(s \cdot n\) equals \((s|_V) \cdot (n|_V)\), so \(a_R\) commutes with the
  triple-tensor and \(\mathbb{Z}\)-tensor restriction maps. The verification
  reduces to the \(\mathcal{O}_X\)-semilinearity of the restriction maps of
  \(\mathcal{Q}\), checked on elementary tensors by \(\otimes\)-induction.
\end{definition}

\begin{definition}

\leanok
[Projection natural transformation of the coequalizer]
  \label{def:relTensorProj}
  \lean{AlgebraicGeometry.Scheme.Modules.relTensorProj}
  \uses{def:relTensorDomainPresheaf, def:snap_objRestrict}
  For presheaves of \(\mathcal{O}_X\)-modules \(\mathcal{P}, \mathcal{Q}\) the
  \emph{projection} natural transformation
  \[
    \pi : \bigl(U \mapsto \mathcal{P}(U) \otimes_{\mathbb{Z}} \mathcal{Q}(U)\bigr)
      \;\longrightarrow\;
      \bigl(U \mapsto \mathcal{P}(U) \otimes_{\mathcal{O}_X(U)} \mathcal{Q}(U)\bigr)
  \]
  from the \(\mathbb{Z}\)-tensor presheaf (\cref{def:relTensorDomainPresheaf}) to
  the underlying abelian-group presheaf of the presheaf-level monoidal tensor
  \(\mathcal{P} \otimes_p \mathcal{Q}\) (Mathlib's
  \(\mathrm{PresheafOfModules.Monoidal.tensorObj}\)) is, on an open \(U\), the
  canonical quotient additive map
  \[
    \mathcal{P}(U) \otimes_{\mathbb{Z}} \mathcal{Q}(U)
    \longrightarrow
    \mathcal{P}(U) \otimes_{\mathcal{O}_X(U)} \mathcal{Q}(U),
    \qquad
    m \otimes n \longmapsto m \otimes_{\mathcal{O}_X(U)} n,
  \]
  the objectwise projection \(\mathrm{projL}\) of
  \cref{lem:relativeTensor_objectwise_coequalizer}. It is the cofork of the
  relative-tensor coequalizer: the third leg, the cokernel of \(a_L - a_R\)
  (\cref{def:relTensorActL}, \cref{def:relTensorActR}). The apex is identified
  with the objectwise value of the presheaf monoidal product through
  \cref{lem:presheaf_tensorObj_obj_mathlib}.

  \emph{Naturality.} The objectwise square asserts
  \[
    \pi_V \circ
      \bigl(\mathrm{restr}_{\mathcal{P}} \otimes_{\mathbb{Z}}
        \mathrm{restr}_{\mathcal{Q}}\bigr)
    \;=\;
    \bigl((\mathcal{P} \otimes_p \mathcal{Q})(f)\bigr) \circ \pi_U
    \qquad (f : V \to U),
  \]
  both composites sending \(m \otimes n\) to
  \(\bigl(m|_V\bigr) \otimes_{\mathcal{O}_X(V)} \bigl(n|_V\bigr)\). On elementary
  tensors the left-hand side reduces by functoriality of the
  \(\mathbb{Z}\)-tensor of restriction maps followed by the defining formula of
  \(\mathrm{projL}\); the right-hand side reduces by \(\mathrm{projL}\) followed
  by the objectwise restriction formula for the presheaf monoidal product
  (\cref{lem:presheaf_tensorObj_obj_mathlib}, on elementary tensors:
  \((\mathcal{P} \otimes_p \mathcal{Q})(f)\,(m \otimes_{\mathcal{O}_X(U)} n) =
  (m|_V) \otimes_{\mathcal{O}_X(V)} (n|_V)\)). The two sides then agree by
  \(\otimes\)-induction.

\end{definition}

\begin{lemma}

\leanok
[Cofork condition of the relative-tensor coequalizer]
  \label{lem:snap_relTensorActL_proj_eq}
  \lean{AlgebraicGeometry.Scheme.Modules.relTensorActL_proj_eq}
  \uses{def:relTensorActL, def:relTensorActR, def:relTensorProj,
    lem:relativeTensor_objectwise_coequalizer}
  For presheaves of \(\mathcal{O}_X\)-modules \(\mathcal{P}, \mathcal{Q}\) the
  left- and right-action transformations compose equally with the projection,
  \[
    \pi \circ a_L \;=\; \pi \circ a_R,
  \]
  as natural transformations from the triple presheaf
  (\cref{def:relTensorTriplePresheaf}) to the underlying abelian-group presheaf of
  \(\mathcal{P} \otimes_p \mathcal{Q}\), where \(a_L\)
  (\cref{def:relTensorActL}), \(a_R\) (\cref{def:relTensorActR}) are the two action
  rows and \(\pi\) (\cref{def:relTensorProj}) is the projection. This is the cofork
  (coequalizing) condition of the presheaf-level relative-tensor coequalizer
  presentation \cref{lem:relativeTensor_as_coequalizer}. It is checked
  componentwise: at each open \(U\) both sides agree by the objectwise coequalizing
  identity \(\pi_U \circ a_{L,U} = \pi_U \circ a_{R,U}\) of
  \cref{lem:relativeTensor_objectwise_coequalizer}.
\end{lemma}

\begin{lemma}\leanok
[Right-whiskering preserves local isomorphisms of presheaves of modules]
  \label{lem:snap_ztensor_whisker_localIso}
  \lean{AlgebraicGeometry.Scheme.Modules.ztensor_whisker_localIso}
  \uses{lem:localIso_toPresheaf_map_unit, def:relTensorDomainPresheaf,
    def:relTensorTriplePresheaf, lem:relativeTensor_as_coequalizer,
    def:snap_opensTopology, lem:presheafModule_monoidal_mathlib}
  Let \(f : \mathcal{P} \to \mathcal{Q}\) be a morphism of presheaves of
  \(\mathcal{O}_X\)-modules whose underlying abelian-presheaf morphism is a
  local isomorphism (lies in the class \(J.W\) for the opens topology \(J\)
  on \(X\)). Then for any presheaf of \(\mathcal{O}_X\)-modules \(\mathcal{R}\),
  the underlying abelian morphism of the right-whiskered map
  \(f \rhd \mathcal{R} : \mathcal{P} \otimes_p \mathcal{R} \to
  \mathcal{Q} \otimes_p \mathcal{R}\) (in the presheaf monoidal structure of
  \cref{lem:presheafModule_monoidal_mathlib}) again lies in \(J.W\).
\end{lemma}

\begin{proof}
  \uses{def:relTensorDomainPresheaf, def:relTensorTriplePresheaf,
    lem:relativeTensor_as_coequalizer, lem:relativeTensor_objectwise_coequalizer}
    \leanok
  Two reductions.

  \emph{(i) \(\mathbb{Z}\)-tensor whiskering preserves \(J.W\).} The category
  of abelian-group presheaves on \(X\) embeds, by a carrier-preserving
  equivalence, into presheaves valued in modules over a (universe lift of)
  \(\mathbb{Z}\). The latter value category is braided monoidal closed, so the
  pointwise (Day) monoidal structure on its presheaf category satisfies: right
  whiskering by a fixed object carries \(J.W\) into \(J.W\) (the reflection
  argument for monoidal sites — sheafification is a monoidal localization for a
  braided closed value category). An equivalence of value categories is a left
  adjoint in both directions, hence preserves sheafification both ways and
  identifies the two classes \(J.W\); transporting along it shows that the
  \(\mathbb{Z}\)-tensor whiskering \(g \otimes_{\mathbb{Z}} \mathrm{id}\) of any
  \(g \in J.W\) by a fixed abelian presheaf again lies in \(J.W\).

  \emph{(ii) Descent through the coequalizer presentation.} By
  \cref{lem:relativeTensor_as_coequalizer}, the underlying abelian presheaf of
  \(\mathcal{P} \otimes_p \mathcal{R}\) is the coequalizer (computed objectwise)
  of the two action rows out of the triple \(\mathbb{Z}\)-tensor presheaf
  (\cref{def:relTensorTriplePresheaf}), and the underlying abelian morphism of
  \(f \rhd \mathcal{R}\) is the map induced on coequalizer points by
  \(\mathbb{Z}\)-whiskering the two rows with \(f\). Abelian sheafification is
  a left adjoint, hence preserves these coequalizers; by (i) it inverts the
  whiskered rows, hence it inverts the induced map on coequalizer points. A
  morphism inverted by sheafification lies in \(J.W\), which is the claim.
\end{proof}

\begin{lemma}[Abelian universal property of the relative tensor product]
  \label{lem:tensorProduct_liftAddHom_mathlib}
  \lean{TensorProduct.liftAddHom}
  \mathlibok
  \textit{Provided by Mathlib (\texttt{Mathlib.LinearAlgebra.TensorProduct.Basic}).}
  Let \(S\) be a commutative ring and \(M, N, P\) be \(S\)-modules (or, more
  generally, the abelian-group statement: \(M, N\) modules and \(P\) an additive
  commutative group). An additive-bilinear map
  \(f : M \to (N \to P)\) that is \emph{\(S\)-balanced}, i.e.\ satisfies
  \(f(s \cdot m)(n) = f(m)(s \cdot n)\) for all \(s \in S\), factors uniquely
  through the relative tensor product as an additive group homomorphism
  \(\widetilde{f} : M \otimes_S N \to P\) with
  \(\widetilde{f}(m \otimes_S n) = f(m)(n)\). This is the abelian-group universal
  property of \(M \otimes_S N\): it is the quotient of \(M \otimes_{\mathbb{Z}} N\)
  by the \(S\)-balancing relations.
\end{lemma}

\begin{lemma}[The relative tensor of modules is a coequalizer of abelian groups]
  \label{lem:relativeTensor_objectwise_coequalizer}
  \lean{AlgebraicGeometry.Scheme.Modules.RelativeTensorCoequalizer.isColimitCofork,
        AlgebraicGeometry.Scheme.Modules.RelativeTensorCoequalizer.actN,
        AlgebraicGeometry.Scheme.Modules.RelativeTensorCoequalizer.actM,
        AlgebraicGeometry.Scheme.Modules.RelativeTensorCoequalizer.actLmap,
        AlgebraicGeometry.Scheme.Modules.RelativeTensorCoequalizer.actRmap,
        AlgebraicGeometry.Scheme.Modules.RelativeTensorCoequalizer.projL,
        AlgebraicGeometry.Scheme.Modules.RelativeTensorCoequalizer.projL_surjective,
        AlgebraicGeometry.Scheme.Modules.RelativeTensorCoequalizer.projL_comp_act,
        AlgebraicGeometry.Scheme.Modules.RelativeTensorCoequalizer.aL,
        AlgebraicGeometry.Scheme.Modules.RelativeTensorCoequalizer.aR,
        AlgebraicGeometry.Scheme.Modules.RelativeTensorCoequalizer.piMor,
        AlgebraicGeometry.Scheme.Modules.RelativeTensorCoequalizer.piMor_epi,
        AlgebraicGeometry.Scheme.Modules.RelativeTensorCoequalizer.coeq_condition,
        AlgebraicGeometry.Scheme.Modules.RelativeTensorCoequalizer.cofork,
        AlgebraicGeometry.Scheme.Modules.RelativeTensorCoequalizer.descHom,
        AlgebraicGeometry.Scheme.Modules.RelativeTensorCoequalizer.descMor,
        AlgebraicGeometry.Scheme.Modules.RelativeTensorCoequalizer.descFac,
        AlgebraicGeometry.Scheme.Modules.RelativeTensorCoequalizer.actLmap_tmul,
        AlgebraicGeometry.Scheme.Modules.RelativeTensorCoequalizer.actRmap_tmul,
        AlgebraicGeometry.Scheme.Modules.RelativeTensorCoequalizer.projL_tmul,
        AlgebraicGeometry.Scheme.Modules.RelativeTensorCoequalizer.piMor_apply,
        AlgebraicGeometry.Scheme.Modules.RelativeTensorCoequalizer.descHom_tmul}
  \uses{lem:tensorProduct_liftAddHom_mathlib}
  Let \(S\) be a commutative ring and \(M, N\) be \(S\)-modules. The abelian group
  \(M \otimes_S N\) is the coequalizer, \emph{in the category
  \(\mathrm{AddCommGrp}\) of abelian groups}, of the two parallel
  \(\mathbb{Z}\)-linear \(S\)-action maps
  \[
    a_L,\, a_R :
    M \otimes_{\mathbb{Z}} (S \otimes_{\mathbb{Z}} N)
    \;\rightrightarrows\;
    M \otimes_{\mathbb{Z}} N,
  \]
  where the \emph{left} action map sends \(m \otimes (s \otimes n) \mapsto
  (s \cdot m) \otimes n\) and the \emph{right} action map sends
  \(m \otimes (s \otimes n) \mapsto m \otimes (s \cdot n)\). The coforking map is
  the canonical projection \(\pi : M \otimes_{\mathbb{Z}} N \to M \otimes_S N\),
  \(m \otimes n \mapsto m \otimes_S n\).
\end{lemma}

\begin{proof}  Write \(a_N : S \otimes_{\mathbb{Z}} N \to N\) and
  \(a_M : M \otimes_{\mathbb{Z}} S \to M\) for the \(\mathbb{Z}\)-linear maps
  induced by the scalar action (\(s \otimes n \mapsto s \cdot n\),
  \(m \otimes s \mapsto s \cdot m\)); the right action map is
  \(a_R = \mathrm{id}_M \otimes a_N\) and the left action map is
  \(a_L = (a_M \otimes \mathrm{id}_N) \circ \mathrm{assoc}^{-1}\), where
  \(\mathrm{assoc}\) is the \(\mathbb{Z}\)-tensor associator. The projection
  \(\pi\) is the additive homomorphism produced by the abelian universal property
  \cref{lem:tensorProduct_liftAddHom_mathlib} from the \(S\)-balanced bilinear map
  \((m, n) \mapsto m \otimes_S n\).

  \emph{Cofork condition.} On an elementary tensor
  \(m \otimes (s \otimes n)\) we have
  \(\pi(a_L(m \otimes s \otimes n)) = (s \cdot m) \otimes_S n
  = m \otimes_S (s \cdot n) = \pi(a_R(m \otimes s \otimes n))\), the middle
  equality being the defining \(S\)-balancing relation of \(M \otimes_S N\); since
  elementary tensors generate, \(\pi \circ a_L = \pi \circ a_R\).

  \emph{Universal property.} Given any abelian group \(P\) and a homomorphism
  \(g : M \otimes_{\mathbb{Z}} N \to P\) with \(g \circ a_L = g \circ a_R\), the
  bilinear map \((m, n) \mapsto g(m \otimes n)\) is \(S\)-balanced (that is exactly
  the cofork condition read on elementary tensors), so by
  \cref{lem:tensorProduct_liftAddHom_mathlib} it factors as a unique
  \(\widetilde{g} : M \otimes_S N \to P\) with \(\widetilde{g} \circ \pi = g\).
  Existence is this descended map; uniqueness follows because \(\pi\) is
  surjective, hence an epimorphism of abelian groups, so \(\widetilde{g} \circ \pi
  = g\) determines \(\widetilde{g}\). Packaging existence and uniqueness gives the
  colimit cofork.

\end{proof}

\begin{lemma}\leanok
[Relative tensor of module presheaves as a coequalizer]
  \label{lem:relativeTensor_as_coequalizer}
  \lean{AlgebraicGeometry.Scheme.Modules.relativeTensorCoequalizerIso}
  \uses{lem:presheafModule_monoidal_mathlib,
    lem:relativeTensor_objectwise_coequalizer,
    lem:evaluationJointlyReflectsColimits_mathlib,
    lem:presheaf_tensorObj_obj_mathlib,
    def:relTensorActL, def:relTensorActR, def:relTensorProj}
  Let \(\mathcal{P}, \mathcal{Q}\) be presheaves of \(\mathcal{O}_X\)-modules. The
  underlying abelian-group presheaf of their relative tensor product
  \(\mathcal{P} \otimes_{p,\mathcal{O}_X} \mathcal{Q}\)
  (\cref{lem:presheafModule_monoidal_mathlib}) is naturally the coequalizer of the
  parallel pair
  \[
    \mathcal{P} \otimes_{\mathbb{Z}} \mathcal{O}_X \otimes_{\mathbb{Z}} \mathcal{Q}
    \;\rightrightarrows\;
    \mathcal{P} \otimes_{\mathbb{Z}} \mathcal{Q}
  \]
  of abelian-group presheaves, where the first map acts the ring factor on the left,
  \(p \otimes r \otimes q \mapsto (r \cdot p) \otimes q\), and the second on the
  right, \(p \otimes r \otimes q \mapsto p \otimes (r \cdot q)\). The identification
  is objectwise and natural in \(U\).
\end{lemma}

\begin{proof}
  \uses{lem:presheafModule_monoidal_mathlib,
    lem:relativeTensor_objectwise_coequalizer,
    lem:evaluationJointlyReflectsColimits_mathlib,
    lem:presheaf_tensorObj_obj_mathlib,
    def:relTensorActL, def:relTensorActR, def:relTensorProj,
    lem:snap_relTensorActL_proj_eq}
    \leanok
  The argument promotes the objectwise coequalizer of
  \cref{lem:relativeTensor_objectwise_coequalizer} to the functor category, in
  three steps.

  \textbf{1. Objectwise.} For each open \(U\), apply
  \cref{lem:relativeTensor_objectwise_coequalizer} with \(S = \mathcal{O}_X(U)\),
  \(M = \mathcal{P}(U)\), \(N = \mathcal{Q}(U)\): the abelian group
  \(\mathcal{P}(U) \otimes_{\mathcal{O}_X(U)} \mathcal{Q}(U)\) is the coequalizer
  in \(\mathrm{AddCommGrp}\) of the two \(\mathcal{O}_X(U)\)-action maps
  \[
    \mathcal{P}(U) \otimes_{\mathbb{Z}} \mathcal{O}_X(U)
      \otimes_{\mathbb{Z}} \mathcal{Q}(U)
    \;\rightrightarrows\;
    \mathcal{P}(U) \otimes_{\mathbb{Z}} \mathcal{Q}(U),
  \]
  with cofork map the projection \(\pi_U\).

  \textbf{2. Naturality / promotion.} The three objectwise families are precisely
  the natural transformations of presheaves on \(X\)
  \[
    \mathcal{P} \otimes_{\mathbb{Z}} \mathcal{O}_X
      \otimes_{\mathbb{Z}} \mathcal{Q}
    \;\overset{a_L,\,a_R}{\rightrightarrows}\;
    \mathcal{P} \otimes_{\mathbb{Z}} \mathcal{Q}
    \;\overset{\pi}{\longrightarrow}\;
    \bigl(U \mapsto \mathcal{P}(U) \otimes_{\mathcal{O}_X(U)} \mathcal{Q}(U)\bigr)
  \]
  built in the previous definitions: the parallel pair is the left- and
  right-action transformations \(a_L\) (\cref{def:relTensorActL}) and \(a_R\)
  (\cref{def:relTensorActR}) from the triple presheaf
  (\cref{def:relTensorTriplePresheaf}) to the \(\mathbb{Z}\)-tensor presheaf
  (\cref{def:relTensorDomainPresheaf}), and the cofork leg is the projection
  \(\pi\) (\cref{def:relTensorProj}) onto the apex. These are morphisms of
  functors \((\mathrm{Opens}\,X)^{\mathrm{op}} \to \mathrm{Ab}\); naturality of
  each in \(U\) is the compatibility of the module action with the restriction
  maps of \(\mathcal{P}, \mathcal{O}_X, \mathcal{Q}\) (the restriction of
  \(s \cdot m\) is \((s|_V) \cdot (m|_V)\), and likewise for the projection),
  recorded in those definitions. Objectwise at each \(U\) the cofork
  \((a_L, a_R, \pi)_U\) is the colimit cofork of step 1. Colimits in a functor
  category are computed objectwise --- by
  \cref{lem:evaluationJointlyReflectsColimits_mathlib} a cocone is a colimit as
  soon as every evaluation \((\mathrm{ev}_U)\) of it is --- so the objectwise
  coequalizers of step 1 promote to a coequalizer of the displayed parallel pair
  \(a_L, a_R\) with cofork \(\pi\) in
  \((\mathrm{Opens}\,X)^{\mathrm{op}} \to \mathrm{Ab}\).

  \textbf{3. Identification of the apex.} The apex presheaf
  \(U \mapsto \mathcal{P}(U) \otimes_{\mathcal{O}_X(U)} \mathcal{Q}(U)\) is, by the
  objectwise tensor formula of the presheaf monoidal structure
  (\cref{lem:presheaf_tensorObj_obj_mathlib}: \((\mathcal{M}_1 \otimes_p
  \mathcal{M}_2)(U) = \mathcal{M}_1(U) \otimes_{\mathcal{O}_X(U)}
  \mathcal{M}_2(U)\)), exactly the underlying abelian-group presheaf of the
  relative tensor product \(\mathcal{P} \otimes_{p,\mathcal{O}_X} \mathcal{Q}\)
  (\cref{lem:presheafModule_monoidal_mathlib}). Transporting the colimit of step 2
  along this identification exhibits
  \(\mathcal{P} \otimes_{p,\mathcal{O}_X} \mathcal{Q}\), as an abelian-group
  presheaf, as the coequalizer of the parallel pair, naturally in \(U\).
\end{proof}

\begin{lemma}[Objectwise colimits in a functor category]
  \label{lem:evaluationJointlyReflectsColimits_mathlib}
  \lean{CategoryTheory.Limits.evaluationJointlyReflectsColimits}
  \mathlibok
  \textit{Provided by Mathlib
  (\texttt{Mathlib.CategoryTheory.Limits.FunctorCategory.Basic}).}
  Let \(\mathcal{D}\) be a category and \(\mathcal{A}\) a category in which the
  relevant colimits exist. A cocone over a diagram
  \(F : \mathcal{J} \to (\mathcal{D} \to \mathcal{A})\) in the functor category is
  a colimit cocone as soon as, for every object \(U\) of \(\mathcal{D}\), its image
  under evaluation \(\mathrm{ev}_U : (\mathcal{D} \to \mathcal{A}) \to \mathcal{A}\)
  is a colimit cocone in \(\mathcal{A}\). Equivalently, colimits in a functor
  category are computed objectwise.
\end{lemma}

\begin{lemma}[Objectwise value of the presheaf tensor product]
  \label{lem:presheaf_tensorObj_obj_mathlib}
  \lean{PresheafOfModules.Monoidal.tensorObj_obj}
  \mathlibok
  \textit{Provided by Mathlib
  (\texttt{Mathlib.Algebra.Category.ModuleCat.Presheaf.Monoidal}).}
  Let \(R : \mathcal{C}^{\mathrm{op}} \to \mathrm{CommRingCat}\) be a presheaf of
  commutative rings and \(\mathcal{M}_1, \mathcal{M}_2\) presheaves of modules over
  \(R\). For every object \(U\) the value of their monoidal product
  (\cref{lem:presheafModule_monoidal_mathlib}) is the objectwise relative tensor
  product of modules,
  \[
    (\mathcal{M}_1 \otimes_p \mathcal{M}_2)(U)
      \;=\; \mathcal{M}_1(U) \otimes_{R(U)} \mathcal{M}_2(U),
  \]
  with the restriction maps the tensor products of the restriction maps.
\end{lemma}

\begin{lemma}\leanok
[Sheafification inverts the whiskered localization unit]
  \label{lem:isIso_sheafification_whiskerRight_unit}
  \lean{AlgebraicGeometry.Scheme.Modules.isIso_sheafification_whiskerRight_unit}
  \uses{lem:presheafModule_monoidal_mathlib,
    lem:presheafModule_sheafification_mathlib,
    lem:isIso_sheafification_map_iff, lem:localIso_toPresheaf_map_unit,
    lem:relativeTensor_as_coequalizer,
    lem:relativeTensor_objectwise_coequalizer}
  Let \(\mathcal{P}, \mathcal{Q}\) be presheaves of \(\mathcal{O}_X\)-modules,
  let \(\eta_{\mathcal{P}} : \mathcal{P} \to (\mathcal{P})^{\#}\) be the unit of
  the sheafification adjunction
  (\cref{lem:presheafModule_sheafification_mathlib}), viewed as a morphism of
  presheaves of modules, and write
  \[
    \eta_{\mathcal{P}} \rhd \mathcal{Q} :
    \mathcal{P} \otimes_p \mathcal{Q}
    \longrightarrow (\mathcal{P})^{\#} \otimes_p \mathcal{Q}
  \]
  for its right whiskering by \(\mathcal{Q}\) in the presheaf monoidal structure
  (\cref{lem:presheafModule_monoidal_mathlib}). Then the image of
  \(\eta_{\mathcal{P}} \rhd \mathcal{Q}\) under the sheafification functor,
  \[
    (\eta_{\mathcal{P}} \rhd \mathcal{Q})^{\#} :
    (\mathcal{P} \otimes_p \mathcal{Q})^{\#}
    \xrightarrow{\ \sim\ }
    \bigl((\mathcal{P})^{\#} \otimes_p \mathcal{Q}\bigr)^{\#},
  \]
  is an isomorphism of sheaves of modules. This is the strong-monoidality
  comparison of the module sheafification functor on a whiskered unit.
\end{lemma}

\begin{proof}
  \uses{lem:presheafModule_monoidal_mathlib,
    lem:presheafModule_sheafification_mathlib,
    lem:isIso_sheafification_map_iff, lem:localIso_toPresheaf_map_unit,
    lem:relativeTensor_as_coequalizer,
    lem:relativeTensor_objectwise_coequalizer,
    lem:snap_ztensor_whisker_localIso, def:relTensorDomainPresheaf}
    \leanok
  Module sheafification is a localization functor: it inverts exactly the class
  \(W\) of morphisms of presheaves of modules whose underlying morphism of
  abelian-group presheaves is a local isomorphism --- that is, lies in the class
  \(J.W\) of weak equivalences attached to the Grothendieck topology \(J\) on
  \(X\) (the inverse image of \(J.W\) along the forgetful functor to
  abelian-group presheaves). By the localization criterion for a reflective
  adjunction, \((f)^{\#}\) is an isomorphism if and only if \(f \in W\), i.e.\
  if and only if the underlying abelian-presheaf morphism of \(f\) lies in
  \(J.W\). It therefore suffices to prove that the underlying abelian morphism
  of \(\eta_{\mathcal{P}} \rhd \mathcal{Q}\) lies in \(J.W\).

  \textbf{The unit is a local isomorphism.} The sheafification unit
  \(\eta_{\mathcal{P}} : \mathcal{P} \to (\mathcal{P})^{\#}\) has underlying
  abelian-presheaf morphism in \(J.W\)
  (\cref{lem:localIso_toPresheaf_map_unit} — the defining property of the
  associated-sheaf unit for the topology \(J\)).

  \textbf{Whiskering preserves \(J.W\).} By
  \cref{lem:snap_ztensor_whisker_localIso}, right-whiskering a morphism of
  presheaves of modules whose underlying abelian morphism lies in \(J.W\) by
  the fixed presheaf \(\mathcal{Q}\) yields a morphism whose underlying
  abelian morphism again lies in \(J.W\) (proved there by descending the
  pointwise monoidal-whiskering statement through the relative-tensor
  coequalizer presentation of
  \cref{lem:relativeTensor_as_coequalizer}). Applied to
  \(\eta_{\mathcal{P}}\), the underlying abelian morphism of
  \(\eta_{\mathcal{P}} \rhd \mathcal{Q}\) lies in \(J.W\), so
  \(\eta_{\mathcal{P}} \rhd \mathcal{Q} \in W\), and
  \((\eta_{\mathcal{P}} \rhd \mathcal{Q})^{\#}\) is an isomorphism.
\end{proof}

\begin{corollary}\leanok
[Sheaf-level associator on tensor products]
  \label{cor:sheafTensorObjAssoc}
  \lean{AlgebraicGeometry.Scheme.Modules.tensorObjAssoc}
  \uses{def:sheafTensorObj, lem:isIso_sheafification_whiskerRight_unit,
    lem:presheafModule_monoidal_mathlib,
    lem:presheafModule_sheafification_mathlib}
  For sheaves of \(\mathcal{O}_X\)-modules \(\mathcal{A}, \mathcal{B},
  \mathcal{C}\) there is a canonical associativity isomorphism
  \[
    \alpha_{\mathcal{A},\mathcal{B},\mathcal{C}} :
    (\mathcal{A} \otimes_{\mathcal{O}_X} \mathcal{B})
      \otimes_{\mathcal{O}_X} \mathcal{C}
    \;\xrightarrow{\ \sim\ }\;
    \mathcal{A} \otimes_{\mathcal{O}_X}
      (\mathcal{B} \otimes_{\mathcal{O}_X} \mathcal{C})
  \]
  of sheaves of modules (\cref{def:sheafTensorObj}), natural in all three
  arguments.
\end{corollary}

\begin{proof}
  \uses{def:sheafTensorObj, lem:isIso_sheafification_whiskerRight_unit,
    lem:presheafModule_monoidal_mathlib,
    lem:presheafModule_sheafification_mathlib}
    \leanok
  Write \(a, b, c\) for the underlying presheaves of modules of
  \(\mathcal{A}, \mathcal{B}, \mathcal{C}\) and \((-)^{\#}\) for sheafification
  (\cref{lem:presheafModule_sheafification_mathlib}). By
  \cref{def:sheafTensorObj} the two iterated products unfold to
  \[
    (\mathcal{A} \otimes \mathcal{B}) \otimes \mathcal{C}
      = \bigl((a \otimes_p b)^{\#} \otimes_p c\bigr)^{\#},
    \qquad
    \mathcal{A} \otimes (\mathcal{B} \otimes \mathcal{C})
      = \bigl(a \otimes_p (b \otimes_p c)^{\#}\bigr)^{\#},
  \]
  each carrying an \emph{inner} sheafification that obstructs a direct appeal to
  the presheaf associator. The whiskered units remove these inner
  sheafifications. Applying \cref{lem:isIso_sheafification_whiskerRight_unit}
  with \(\mathcal{P} = a \otimes_p b\) and \(\mathcal{Q} = c\), the morphism
  \[
    (\eta_{a \otimes_p b} \rhd c)^{\#} :
    \bigl((a \otimes_p b) \otimes_p c\bigr)^{\#}
    \xrightarrow{\ \sim\ }
    \bigl((a \otimes_p b)^{\#} \otimes_p c\bigr)^{\#}
    = (\mathcal{A} \otimes \mathcal{B}) \otimes \mathcal{C}
  \]
  is an isomorphism; and, reading the same lemma through the braiding of the
  presheaf monoidal structure (\cref{lem:presheafModule_monoidal_mathlib}) so as
  to whisker the unit \(\eta_{b \otimes_p c}\) on the left of \(a\), the
  morphism
  \[
    (a \lhd \eta_{b \otimes_p c})^{\#} :
    \bigl(a \otimes_p (b \otimes_p c)\bigr)^{\#}
    \xrightarrow{\ \sim\ }
    \bigl(a \otimes_p (b \otimes_p c)^{\#}\bigr)^{\#}
    = \mathcal{A} \otimes (\mathcal{B} \otimes \mathcal{C})
  \]
  is an isomorphism. Between the two un-sheafified middle terms sits the
  sheafification of the presheaf associator
  \[
    \alpha^p_{a,b,c} :
    (a \otimes_p b) \otimes_p c
    \xrightarrow{\ \sim\ }
    a \otimes_p (b \otimes_p c)
  \]
  (\cref{lem:presheafModule_monoidal_mathlib}), an isomorphism because
  sheafification is a functor and \(\alpha^p\) is one. The associator is the
  composite
  \[
    \alpha_{\mathcal{A},\mathcal{B},\mathcal{C}}
    = (a \lhd \eta_{b \otimes_p c})^{\#}
      \circ (\alpha^p_{a,b,c})^{\#}
      \circ \bigl((\eta_{a \otimes_p b} \rhd c)^{\#}\bigr)^{-1},
  \]
  an isomorphism as a composite of isomorphisms. Naturality in
  \(\mathcal{A}, \mathcal{B}, \mathcal{C}\) follows from naturality of the unit,
  of the whiskerings, and of the presheaf associator.
\end{proof}

\begin{lemma}\leanok
[Tensor-power comparison isomorphism]
  \label{lem:sheafTensorPow_add}
  \lean{AlgebraicGeometry.Scheme.Modules.tensorPowAdd}
  \uses{def:sheafTensorPow, def:sheafTensorObj}
  % SOURCE: [Stacks Project], "Sheaves of Modules", \S "Invertible modules"
  % (Tag 01CU, following `definition-powers') (read from
  % references/stacks-modules.tex, L4249-L4254).
  % SOURCE QUOTE:
  % "With this definition we have canonical isomorphisms
  %  $\mathcal{L}^{\otimes n} \otimes_{\mathcal{O}_X}
  %  \mathcal{L}^{\otimes m} \to
  %  \mathcal{L}^{\otimes n + m}$, and these isomorphisms
  %  satisfy a commutativity and an associativity constraint
  %  (formulation omitted)."
  \textit{Source: [Stacks Project], "Sheaves of Modules", Tag 01CU (canonical
  power isomorphisms).}
  For a line bundle \(\mathcal{L}\) and \(m, m' \in \mathbb{N}\) there is a
  canonical isomorphism of sheaves of \(\mathcal{O}_X\)-modules
  \[
    \mu_{m,m'} :
    \mathcal{L}^{\otimes m} \otimes_{\mathcal{O}_X} \mathcal{L}^{\otimes m'}
    \;\xrightarrow{\ \sim\ }\; \mathcal{L}^{\otimes (m+m')},
  \]
  natural in \(\mathcal{L}\), and these isomorphisms satisfy a commutativity
  constraint (\(\mu_{m,m'}\) agrees with \(\mu_{m',m}\) after the braiding) and
  an associativity constraint (the two ways of combining
  \(\mathcal{L}^{\otimes m} \otimes \mathcal{L}^{\otimes m'} \otimes
  \mathcal{L}^{\otimes m''}\) into \(\mathcal{L}^{\otimes (m+m'+m'')}\) agree).
\end{lemma}

\begin{proof}
  \uses{def:sheafTensorPow, def:sheafTensorObj, cor:sheafTensorObjAssoc,
    def:tensorObjUnitIso, def:tensorBraiding,
    lem:isIso_sheafification_whiskerRight_unit,
    lem:presheafModule_monoidal_mathlib,
    lem:presheafModule_sheafification_mathlib}
    \leanok
  We do \emph{not} build a full monoidal structure on the whole category
  \(\ShMod(\mathcal{O}_X)\) --- none is constructed here --- but only assemble
  the comparison \(\mu_{m,m'}\) from those structural isomorphisms of
  \cref{def:sheafTensorObj} that the induction actually uses: the associator,
  the braiding, and the left unitor.

  \textbf{The structural maps.} The braiding and the unitors descend to
  \(\ShMod(\mathcal{O}_X)\) directly: each is the sheafification of the
  corresponding isomorphism of the presheaf symmetric monoidal structure
  (\cref{lem:presheafModule_monoidal_mathlib}), and sheafification, being a
  functor, carries isomorphisms to isomorphisms. The \emph{associator} is not
  the sheafification of a single presheaf isomorphism, because the iterated
  sheaf tensor products
  \((\mathcal{A} \otimes \mathcal{B}) \otimes \mathcal{C}\) and
  \(\mathcal{A} \otimes (\mathcal{B} \otimes \mathcal{C})\) each carry an inner
  sheafification (\cref{def:sheafTensorObj}); it is supplied instead by
  \cref{cor:sheafTensorObjAssoc}, whose construction clears those inner
  sheafifications using the strong-monoidality comparison
  \cref{lem:isIso_sheafification_whiskerRight_unit}. These are the only
  isomorphisms of sheaves of modules the construction invokes; right-exactness
  of the tensor product is never used.

  \textbf{Construction of \(\mu_{m,m'}\).} Fix \(m'\) and induct on \(m\). For
  \(m = 0\),
  \[
    \mathcal{L}^{\otimes 0} \otimes \mathcal{L}^{\otimes m'}
      = \mathcal{O}_X \otimes \mathcal{L}^{\otimes m'}
      \xrightarrow{\ \sim\ } \mathcal{L}^{\otimes m'}
  \]
  is the left unitor of \cref{def:sheafTensorObj}. For the inductive step the
  recursive definition
  \(\mathcal{L}^{\otimes (m+1)} = \mathcal{L}^{\otimes m} \otimes \mathcal{L}\)
  (\cref{def:sheafTensorPow}) gives
  \[
    \mathcal{L}^{\otimes (m+1)} \otimes \mathcal{L}^{\otimes m'}
    = (\mathcal{L}^{\otimes m} \otimes \mathcal{L}) \otimes
      \mathcal{L}^{\otimes m'}
    \xrightarrow{\ \alpha\ } \mathcal{L}^{\otimes m} \otimes
      (\mathcal{L} \otimes \mathcal{L}^{\otimes m'})
    \xrightarrow{\ \mathrm{id} \otimes \beta\ } \mathcal{L}^{\otimes m} \otimes
      (\mathcal{L}^{\otimes m'} \otimes \mathcal{L}),
  \]
  using the associator \(\alpha\) (\cref{cor:sheafTensorObjAssoc}) and then the
  braiding \(\beta\) of \cref{def:sheafTensorObj} whiskered into the left
  factor. The inverse associator \(\alpha^{-1}\)
  (\cref{cor:sheafTensorObjAssoc}) and the inductive comparison \(\mu_{m,m'}\)
  whiskered by \(\mathcal{L}\) on the right then carry this to
  \[
    (\mathcal{L}^{\otimes m} \otimes \mathcal{L}^{\otimes m'}) \otimes
      \mathcal{L}
    \xrightarrow{\ \mu_{m,m'} \rhd \mathcal{L}\ }
    \mathcal{L}^{\otimes (m + m')} \otimes \mathcal{L}
    = \mathcal{L}^{\otimes ((m + m') + 1)},
  \]
  the final equality being the recursive definition. The reindexing
  \((m + m') + 1 = (m + 1) + m'\) identifies the target with
  \(\mathcal{L}^{\otimes ((m+1) + m')}\), and the composite is
  \(\mu_{m+1,m'}\), an isomorphism as a composite of isomorphisms. Naturality in
  \(\mathcal{L}\) holds because each constituent comparison is natural.

  \textbf{Commutativity and associativity constraints.} The braiding and
  unitors are the sheafifications of the symmetric-monoidal coherence data of
  \(\PshMod(\mathcal{O}_X)\) (\cref{lem:presheafModule_monoidal_mathlib}), and
  the associator is intertwined with the presheaf associator by the comparison
  isomorphisms of \cref{lem:isIso_sheafification_whiskerRight_unit} through the
  construction of \cref{cor:sheafTensorObjAssoc}. Both the commutativity square
  (\(\mu_{m,m'}\) versus \(\mu_{m',m}\) postcomposed with the braiding) and the
  associativity pentagon (the two bracketings of
  \(\mathcal{L}^{\otimes m} \otimes \mathcal{L}^{\otimes m'} \otimes
  \mathcal{L}^{\otimes m''} \to \mathcal{L}^{\otimes (m+m'+m'')}\)) are therefore
  images, under the sheafification functor together with these comparison
  isomorphisms, of the corresponding coherence diagrams of the presheaf
  symmetric monoidal structure, which commute by Mac\,Lane coherence. Since a
  functor preserves commuting diagrams and the comparison isomorphisms are
  natural, the constraints hold for the family \(\{\mu_{m,m'}\}\). We emphasise
  that this argument uses only the single strong-monoidality comparison
  \cref{lem:isIso_sheafification_whiskerRight_unit} and the resulting associator
  \cref{cor:sheafTensorObjAssoc}, not a full monoidal structure on
  \(\ShMod(\mathcal{O}_X)\).
\end{proof}

\section{Lax-monoidal global sections}
\label{sec:sgr_lax_sections}

The graded multiplication is supplied by the lax-monoidal structure of the
global-sections functor \(\Gamma(X, -) : \ShMod(\mathcal{O}_X) \to
\mathrm{ModuleCat}(\Gamma(X, \mathcal{O}_X))\). It is only \emph{lax} because
global sections do not commute with sheafification: a global section of a
sheafified tensor product is not in general a finite sum of elementary tensors of
global sections.

\begin{definition}\leanok
[Section multiplication]
  \label{def:sectionMul}
  \lean{AlgebraicGeometry.Scheme.Modules.sectionsMul}
  \uses{def:sheafTensorObj, lem:presheafModule_monoidal_mathlib,
    lem:presheafModule_sheafification_mathlib}
  Let \(\mathcal{F}, \mathcal{G}\) be sheaves of \(\mathcal{O}_X\)-modules. The
  \emph{section multiplication} is the \(\Gamma(X,\mathcal{O}_X)\)-bilinear map
  \[
    \Gamma(X, \mathcal{F}) \otimes_{\Gamma(X,\mathcal{O}_X)}
      \Gamma(X, \mathcal{G})
    \;\longrightarrow\;
    \Gamma(X, \mathcal{F} \otimes_{\mathcal{O}_X} \mathcal{G})
  \]
  defined as follows. A pair of global sections \((\sigma, \tau)\) determines,
  by the objectwise formula of \cref{lem:presheafModule_monoidal_mathlib} at the
  top open, the elementary tensor \(\sigma \otimes \tau \in
  (\mathcal{F} \otimes_{p,\mathcal{O}_X} \mathcal{G})(X)\), a global section of
  the tensor product presheaf. Composing with the global-sections component of
  the sheafification unit
  \(\eta : \mathcal{F} \otimes_{p,\mathcal{O}_X} \mathcal{G} \to
  (\mathcal{F} \otimes_{p,\mathcal{O}_X} \mathcal{G})^{\#} =
  \mathcal{F} \otimes_{\mathcal{O}_X} \mathcal{G}\)
  (\cref{lem:presheafModule_sheafification_mathlib},
  \cref{def:sheafTensorObj}) yields a global section
  \(\sigma \cdot \tau \in
  \Gamma(X, \mathcal{F} \otimes_{\mathcal{O}_X} \mathcal{G})\). Bilinearity over
  \(\Gamma(X,\mathcal{O}_X)\) is inherited from bilinearity of the elementary
  tensor and \(\Gamma(X,\mathcal{O}_X)\)-linearity of \(\eta\).
\end{definition}

The graded-monoid axioms below are not raw equalities between sections of two
different tensor powers --- on \(\mathbb{N}\) the indices \((i+j)+k\) and
\(i+(j+k)\), or \(0+n\) and \(n\), are equal but not definitionally so, so the two
sides a priori inhabit distinct section modules. Following
\texttt{Mathlib.LinearAlgebra.TensorPower.Basic}, we phrase each coherence as an
honest equality in a single section module, obtained by transporting one side
along an index-equality isomorphism, and provide a bridge that repackages such a
transported equality as an equality of dependent pairs.

\begin{definition}[Index-equality transport of section components]
  \label{def:sectionsCast}
  \lean{AlgebraicGeometry.Scheme.Modules.sectionsCast}
  \uses{def:sheafTensorPow}
  Let \(\mathcal{L}\) be an invertible sheaf on \(X\) and let \(h : i = j\) be an
  equality of natural numbers. Applying the global-sections functor
  \(\Gamma(X,-)\) to the canonical isomorphism
  \(\mathcal{L}^{\otimes i} \xrightarrow{\sim} \mathcal{L}^{\otimes j}\) induced by
  \(h\) under the tensor-power functor (\cref{def:sheafTensorPow}) yields the
  \(\Gamma(X,\mathcal{O}_X)\)-linear isomorphism
  \[
    \mathrm{sectionsCast}(h) :
    \Gamma(X, \mathcal{L}^{\otimes i}) \;\xrightarrow{\ \sim\ }\;
    \Gamma(X, \mathcal{L}^{\otimes j}),
  \]
  the \emph{section-component transport} along \(h\). It is the section-level
  analogue of the tensor-power index cast \(\mathrm{TensorPower.cast}\) of
  \texttt{Mathlib.LinearAlgebra.TensorPower.Basic}, which transports
  \(\bigotimes^{i} M\) to \(\bigotimes^{j} M\) along \(i = j\). Because the
  isomorphism induced by the reflexive equality is the identity and
  \(\Gamma(X,-)\) preserves identities, the transport along a reflexive index
  equality is the identity map (the reflexivity simplification
  \(\mathrm{sectionsCast\_refl}\), \cref{lem:sectionsCast_refl}).
\end{definition}

\begin{lemma}[Reflexivity of the section transport]
  \label{lem:sectionsCast_refl}
  \lean{AlgebraicGeometry.Scheme.Modules.sectionsCast_refl}
  \uses{def:sectionsCast}
  For any \(i\), the transport along the reflexive equality \(i = i\)
  (\cref{def:sectionsCast}) equals the identity automorphism of
  \(\Gamma(X, \mathcal{L}^{\otimes i})\). This is the section-level counterpart of
  \(\mathrm{TensorPower.cast\_refl}\), and serves as a simplification lemma
  collapsing the transport along a reflexive index equality.
\end{lemma}

\begin{proof}
  The canonical isomorphism
  \(\mathcal{L}^{\otimes i} \cong \mathcal{L}^{\otimes i}\) induced by the
  reflexive equality is the identity, and the global-sections functor carries
  identities to identities.
\end{proof}

\begin{lemma}[Cast-mediated equality in the graded monoid]
  \label{lem:gradedMonoid_eq_of_cast}
  \lean{AlgebraicGeometry.Scheme.Modules.gradedMonoid_eq_of_cast}
  \uses{def:sectionsCast, lem:sectionsCast_refl}
  Write the carrier of the graded monoid as the dependent sum
  \(\coprod_{n} \Gamma(X, \mathcal{L}^{\otimes n})\), whose elements are pairs
  \((n, s)\) with \(s \in \Gamma(X, \mathcal{L}^{\otimes n})\). Let
  \(a = (i, x)\) and \(b = (j, y)\) be two such pairs. If \(h : i = j\) and the
  transported component agrees, \(\mathrm{sectionsCast}(h)\,x = y\), then
  \(a = b\). That is, a component-level equality mediated by the index transport
  \cref{def:sectionsCast} repackages into a genuine equality of dependent pairs.
  This is the section-level analogue of the bridge
  \(\mathrm{gradedMonoid\_eq\_of\_cast}\) of
  \texttt{Mathlib.LinearAlgebra.TensorPower.Basic} (line~123 there).
\end{lemma}

\begin{proof}
  \uses{lem:sectionsCast_refl}
  Substituting \(j = i\) by means of \(h\) reduces the transport to the identity
  (\cref{lem:sectionsCast_refl}), so the hypothesis becomes \(x = y\); the two
  dependent pairs \((i, x)\) and \((i, y)\) then coincide.
\end{proof}

\begin{lemma}[Coherence of the section multiplication]
  \label{lem:sectionMul_coherent}
  \lean{AlgebraicGeometry.Scheme.Modules.sectionsMul_one_mul,
    AlgebraicGeometry.Scheme.Modules.sectionsMul_mul_one,
    AlgebraicGeometry.Scheme.Modules.sectionsMul_mul_assoc,
    AlgebraicGeometry.Scheme.Modules.sectionsMul_mul_comm}
  \uses{def:sectionMul, def:sectionsCast, lem:moduleUnit_mathlib,
    lem:sheafTensorPow_add}
  Write \(a \cdot b\) for the degreewise section multiplication on tensor powers:
  for \(a \in \Gamma(X, \mathcal{L}^{\otimes i})\) and
  \(b \in \Gamma(X, \mathcal{L}^{\otimes j})\) it is the elementary section
  multiplication (\cref{def:sectionMul}) followed by the tensor-power comparison
  isomorphism \(\mu_{i,j}\) (\cref{lem:sheafTensorPow_add}), landing in
  \(\Gamma(X, \mathcal{L}^{\otimes (i+j)})\); and let \(1\) denote the unit of
  degree \(0\), the image of \(1 \in \Gamma(X,\mathcal{O}_X)\) under the
  identification \(\Gamma(X,\mathcal{O}_X) = \Gamma(X, \mathcal{L}^{\otimes 0})\)
  (\cref{lem:moduleUnit_mathlib}). Mirroring
  \(\mathrm{TensorPower.\{one\_mul, mul\_one, mul\_assoc\}}\) together with the
  commutativity lemma of \texttt{Mathlib.LinearAlgebra.TensorPower.Basic}, the
  multiplication satisfies the following four component-level identities. Each is
  an honest equality in a single section module, obtained by transporting one side
  along the relevant index equality via \cref{def:sectionsCast}; the subscript on
  \(\mathrm{sectionsCast}\) records that index equality.
  \begin{itemize}
    \item \emph{Left unitality.} For \(a \in \Gamma(X, \mathcal{L}^{\otimes n})\),
      \[
        \mathrm{sectionsCast}_{\,0+n=n}\,(1 \cdot a) = a .
      \]
    \item \emph{Right unitality.} For \(a \in \Gamma(X, \mathcal{L}^{\otimes n})\),
      \[
        \mathrm{sectionsCast}_{\,n+0=n}\,(a \cdot 1) = a .
      \]
    \item \emph{Associativity.} For
      \(a \in \Gamma(X, \mathcal{L}^{\otimes n_a})\),
      \(b \in \Gamma(X, \mathcal{L}^{\otimes n_b})\),
      \(c \in \Gamma(X, \mathcal{L}^{\otimes n_c})\),
      \[
        \mathrm{sectionsCast}_{\,(n_a+n_b)+n_c = n_a+(n_b+n_c)}\,
          ((a \cdot b) \cdot c) = a \cdot (b \cdot c) .
      \]
    \item \emph{Commutativity.} For
      \(a \in \Gamma(X, \mathcal{L}^{\otimes n_a})\) and
      \(b \in \Gamma(X, \mathcal{L}^{\otimes n_b})\),
      \[
        \mathrm{sectionsCast}_{\,n_a+n_b = n_b+n_a}\,(a \cdot b) = b \cdot a .
      \]
  \end{itemize}
  The transport is indispensable: the indices on the two sides of each identity
  are equal but not definitionally so, hence a priori live in distinct section
  modules, and \cref{def:sectionsCast} moves one side into the other's module to
  produce a genuine equality there.
\end{lemma}

\begin{proof}
  \uses{def:sectionMul, def:sectionsCast, lem:presheafModule_monoidal_mathlib,
    lem:presheafModule_sheafification_mathlib, lem:sheafTensorPow_add,
    cor:sheafTensorObjAssoc}
  All four identities reduce to the corresponding identities at the level of the
  tensor product presheaf, where by \cref{lem:presheafModule_monoidal_mathlib}
  evaluation at the top open is strictly monoidal, so the associator, unitors and
  symmetry act on elementary tensors by the usual formulas
  \((\sigma \otimes \tau) \otimes \upsilon \mapsto
  \sigma \otimes (\tau \otimes \upsilon)\),
  \(1 \otimes \sigma \mapsto \sigma\) and
  \(\sigma \otimes \tau \mapsto \tau \otimes \sigma\). The sheafification unit
  \(\eta\) is a natural transformation, so it commutes with the associator,
  unitors and symmetry (\cref{lem:presheafModule_sheafification_mathlib}); applying
  \(\Gamma(X,-)\) to the naturality squares transports these presheaf-level
  identities to the section multiplications, and the index transport
  \cref{def:sectionsCast} is precisely the image under \(\Gamma(X,-)\) of the
  tensor-power comparison \(\mu\) (\cref{lem:sheafTensorPow_add}) along which the
  two sides are matched.
\end{proof}

\section{Graded-ring and graded-module assembly}
\label{sec:sgr_graded_assembly}

The components and multiplications of the previous two sections are assembled
into the graded ring and graded module through Mathlib's
\emph{external-direct-sum} graded API: the relevant inputs are families of
abelian groups indexed by \(\mathbb{N}\) together with multiplication maps
landing in the index-sum component, not submodules of a pre-existing ring.

\begin{lemma}[External-direct-sum graded commutative semiring]
  \label{lem:directSum_gcommSemiring_mathlib}
  \lean{DirectSum.GCommSemiring}
  \mathlibok
  \textit{Provided by Mathlib (\texttt{Mathlib.Algebra.DirectSum.Ring}).}
  Let \(A : \mathbb{N} \to \mathrm{Type}\) be a family of additive commutative
  monoids equipped with a graded multiplication
  \(A_i \times A_j \to A_{i+j}\) and a unit \(1 \in A_0\) satisfying graded
  associativity, two-sided unitality, distributivity, and graded commutativity
  (the data of a \(\mathrm{GCommSemiring}\) on \(A\)). Then the external direct
  sum \(\bigoplus_i A_i\) is a commutative semiring, with multiplication the
  bilinear extension of the graded multiplication and identity the image of
  \(1 \in A_0\).
\end{lemma}

\begin{lemma}[External-direct-sum graded module]
  \label{lem:directSum_gmodule_mathlib}
  \lean{DirectSum.Gmodule}
  \mathlibok
  \textit{Provided by Mathlib (\texttt{Mathlib.Algebra.Module.GradedModule}).}
  Let \(A : \mathbb{N} \to \mathrm{Type}\) carry a graded semiring structure as
  in \cref{lem:directSum_gcommSemiring_mathlib} and let
  \(M : \mathbb{N} \to \mathrm{Type}\) be a family of additive commutative
  monoids with a graded scalar action \(A_i \times M_j \to M_{i+j}\) satisfying
  the graded module axioms (the data of a \(\mathrm{Gmodule}\) of \(M\) over
  \(A\)). Then the external direct sum \(\bigoplus_j M_j\) is a module over the
  semiring \(\bigoplus_i A_i\), with scalar multiplication the bilinear
  extension of the graded action.
\end{lemma}

\begin{lemma}[Graded semiring structure on the section components]
  \label{lem:sectionGradedRing_gcommSemiring}
  \lean{AlgebraicGeometry.sectionGradedRing_gcommSemiring}
  \uses{def:sheafTensorPow, lem:sheafTensorPow_add, def:sectionMul,
    lem:sectionMul_coherent, def:sectionsCast, lem:gradedMonoid_eq_of_cast,
    lem:directSum_gcommSemiring_mathlib, lem:moduleUnit_mathlib}
  % SOURCE: [Stacks Project], "Sheaves of Modules", \S "Invertible modules",
  % preceding `definition-gamma-star' (Tag 01CV) (read from
  % references/stacks-modules.tex, L4257-L4267).
  % SOURCE QUOTE:
  % "Let $(X, \mathcal{O}_X)$ be a ringed space. We can define a $\mathbf{Z}$-graded
  %  ring structure on $\bigoplus \Gamma(X, \mathcal{L}^{\otimes n})$ by mapping
  %  $s \in \Gamma(X, \mathcal{L}^{\otimes n})$ and
  %  $t \in \Gamma(X, \mathcal{L}^{\otimes m})$ to the
  %  section corresponding to $s \otimes t$ in
  %  $\Gamma(X, \mathcal{L}^{\otimes n + m})$.
  %  We omit the verification that this defines a commutative
  %  and associative ring with $1$."
  \textit{Source: [Stacks Project], "Sheaves of Modules", Tag 01CV (the graded
  ring structure on \(\bigoplus \Gamma(X,\mathcal{L}^{\otimes n})\)).}
  Let \(L_s\) be a line bundle on \(X_s\). The family of
  \(\Gamma(X_s,\mathcal{O}_{X_s})\)-modules
  \(m \mapsto \Gamma(X_s, L_s^{\otimes m})\) carries a graded commutative
  semiring structure (a \(\mathrm{GCommSemiring}\),
  \cref{lem:directSum_gcommSemiring_mathlib}) whose degree-\((m,m')\)
  multiplication is the composite
  \[
    \Gamma(X_s, L_s^{\otimes m}) \otimes_{\Gamma(X_s,\mathcal{O}_{X_s})}
      \Gamma(X_s, L_s^{\otimes m'})
    \xrightarrow{\ \cdot\ }
    \Gamma\bigl(X_s, L_s^{\otimes m} \otimes L_s^{\otimes m'}\bigr)
    \xrightarrow{\ \Gamma(\mu_{m,m'})\ }
    \Gamma\bigl(X_s, L_s^{\otimes (m+m')}\bigr),
  \]
  the section multiplication (\cref{def:sectionMul}) followed by the
  tensor-power comparison isomorphism (\cref{lem:sheafTensorPow_add}), with unit
  \(1 \in \Gamma(X_s, \mathcal{O}_{X_s}) = \Gamma(X_s, L_s^{\otimes 0})\)
  (\cref{lem:moduleUnit_mathlib}). Consequently
  \(R(X_s, L_s) = \bigoplus_{m \ge 0} \Gamma(X_s, L_s^{\otimes m})\) is a
  commutative semiring; this is the graded-ring structure underlying
  \cref{def:sectionGradedRing}.
\end{lemma}

\begin{proof}
  \uses{lem:sheafTensorPow_add, lem:sectionMul_coherent, def:sectionsCast,
    lem:gradedMonoid_eq_of_cast, def:sectionMul,
    lem:directSum_gcommSemiring_mathlib}
  The structure is assembled field-for-field as \(\mathrm{TensorPower}\) builds its
  graded semiring in \texttt{Mathlib.LinearAlgebra.TensorPower.Basic}. The
  degreewise multiplication \(\Gamma(\mu_{m,m'}) \circ (\,\cdot\,)\) of the
  statement and the degree-\(0\) unit \(1 \in \Gamma(X, \mathcal{L}^{\otimes 0})\)
  provide the graded multiplication and unit. Their graded-monoid axioms --- left
  and right unitality and associativity --- are obtained from the corresponding
  component-level identities of \cref{lem:sectionMul_coherent} by passing each
  through the bridge \cref{lem:gradedMonoid_eq_of_cast}, which converts a
  transport-mediated equality (\cref{def:sectionsCast}) into the required equality
  of dependent pairs; the graded power operation is the default iterated product
  and contributes no further data. This yields the graded-monoid layer.

  The remaining semiring axioms are the bilinearity (left and right distributivity
  and annihilation by \(0\)) of the graded multiplication and the behaviour of the
  natural-number coercion. Bilinearity is immediate: the degreewise multiplication
  \(\Gamma(\mu_{m,m'}) \circ (\,\cdot\,)\) is \(\Gamma(X,\mathcal{O}_X)\)-bilinear,
  being the composite of the bilinear section multiplication (\cref{def:sectionMul})
  with the linear comparison \(\Gamma(\mu_{m,m'})\)
  (\cref{lem:sheafTensorPow_add}), so each side vanishes on \(0\) and distributes
  over sums. The natural-number coercion sends \(n\) to \(n \cdot 1\), the
  \(n\)-fold sum of the degree-\(0\) unit. This completes the graded-semiring
  layer.

  Finally graded commutativity is the commutativity identity of
  \cref{lem:sectionMul_coherent}, again passed through
  \cref{lem:gradedMonoid_eq_of_cast}. These are precisely the data of a
  \(\mathrm{GCommSemiring}\) on \(m \mapsto \Gamma(X, \mathcal{L}^{\otimes m})\),
  so \cref{lem:directSum_gcommSemiring_mathlib} endows
  \(\bigoplus_{m \ge 0} \Gamma(X_s, L_s^{\otimes m})\) with its commutative-semiring
  structure.
\end{proof}

\begin{lemma}[Graded module structure on the twisted-section components]
  \label{lem:sectionGradedModule_gmodule}
  \lean{AlgebraicGeometry.sectionGradedModule_gmodule}
  \uses{def:sheafModuleTwist, lem:sheafTensorPow_add, def:sectionMul,
    lem:sectionMul_coherent, def:sectionsCast, lem:gradedMonoid_eq_of_cast,
    lem:directSum_gmodule_mathlib, lem:sectionGradedRing_gcommSemiring}
  % SOURCE: [Stacks Project], "Sheaves of Modules", \S "Invertible modules",
  % `definition-gamma-star' (Tag 01CV) (read from
  % references/stacks-modules.tex, L4279-L4286).
  % SOURCE QUOTE:
  % "Given a sheaf of $\mathcal{O}_X$-modules $\mathcal{F}$ we set
  %  $$
  %  \Gamma_*(X, \mathcal{L}, \mathcal{F})
  %  =
  %  \bigoplus\nolimits_{n \in \mathbf{Z}} \Gamma(X,
  %  \mathcal{F} \otimes_{\mathcal{O}_X} \mathcal{L}^{\otimes n})
  %  $$
  %  which we think of as a graded $\Gamma_*(X, \mathcal{L})$-module."
  \textit{Source: [Stacks Project], "Sheaves of Modules", Tag 01CV (the graded
  module \(\Gamma_*(X,\mathcal{L},\mathcal{F})\)).}
  Let \(\mathcal{F}_s\) be a coherent sheaf and \(L_s\) a line bundle on
  \(X_s\). The family
  \(m \mapsto \Gamma(X_s, \mathcal{F}_s \otimes L_s^{\otimes m})\)
  (\cref{def:sheafModuleTwist}) carries a graded module structure (a
  \(\mathrm{Gmodule}\), \cref{lem:directSum_gmodule_mathlib}) over the graded
  semiring of \cref{lem:sectionGradedRing_gcommSemiring}, with degree-\((m,m')\)
  action
  \[
    \Gamma(X_s, L_s^{\otimes m}) \otimes_{\Gamma(X_s,\mathcal{O}_{X_s})}
      \Gamma(X_s, \mathcal{F}_s \otimes L_s^{\otimes m'})
    \xrightarrow{\ \cdot\ }
    \Gamma\bigl(X_s, L_s^{\otimes m} \otimes
      (\mathcal{F}_s \otimes L_s^{\otimes m'})\bigr)
    \xrightarrow{\ \sim\ }
    \Gamma\bigl(X_s, \mathcal{F}_s \otimes L_s^{\otimes (m+m')}\bigr),
  \]
  the section multiplication followed by the comparison isomorphism that
  reassociates and merges the line-bundle factors (using
  \cref{lem:sheafTensorPow_add} and the symmetry of the tensor product).
  Consequently \(M(X_s, \mathcal{F}_s, L_s) = \bigoplus_{m \ge 0}
  \Gamma(X_s, \mathcal{F}_s \otimes L_s^{\otimes m})\) is a graded module over
  \(R(X_s, L_s)\); this is the graded-module structure underlying
  \cref{def:sectionGradedModule}.
\end{lemma}

\begin{proof}
  \uses{lem:sectionMul_coherent, lem:sheafTensorPow_add, def:sectionsCast,
    lem:gradedMonoid_eq_of_cast, def:sectionMul,
    lem:directSum_gmodule_mathlib, lem:sectionGradedRing_gcommSemiring}
  The action maps are the displayed composites: the section multiplication
  (\cref{def:sectionMul}) followed by the comparison isomorphism that reassociates
  and merges the line-bundle factors (\cref{lem:sheafTensorPow_add}). The
  \(\mathrm{Gmodule}\) is assembled field-for-field as the graded-action precedent
  \(\mathrm{GradedMonoid.GMulAction}\) of Mathlib, with the action grading on
  indices given by ordinary addition on \(\mathbb{N}\) (so the degree-\((m,m')\)
  action lands in degree \(m+m'\)). Its sigma-type axioms --- unitality
  \(1 \cdot x = x\) and compatibility
  \(r \cdot (r' \cdot x) = (r r') \cdot x\) with the ring multiplication --- are
  the unit and associativity coherences of the section multiplication
  (\cref{lem:sectionMul_coherent}) read with the third tensor factor carrying
  \(\mathcal{F}_s\), each converted from its transport-mediated form
  (\cref{def:sectionsCast}) into the required equality of dependent pairs by the
  bridge \cref{lem:gradedMonoid_eq_of_cast}. The bilinearity axioms of the action
  (additivity and annihilation by \(0\) in each variable) are free, the action map
  being \(\Gamma(X,\mathcal{O}_X)\)-bilinear exactly as in the proof of
  \cref{lem:sectionGradedRing_gcommSemiring}. These are the data of a
  \(\mathrm{Gmodule}\), so \cref{lem:directSum_gmodule_mathlib} produces the module
  structure of \(M(X_s, \mathcal{F}_s, L_s)\) over \(R(X_s, L_s)\).
\end{proof}
